\section*{Outline of Changes}

\begin{itemize}
  \item Added a textual example to the introduction.
  \item Added a TikZ diagram to the introduction to illustrate the motivating example graph and query.
  \item Added a simple TikZ diagram illustrating nodes, edges, labels, and properties for the property graph model in Section~2.1.
  \item Included a table distinguishing between walks, trails, simple paths, and acyclic paths in Section~2.2 to enhance clarity.
  \item Provided short example queries in GQL, SQL/PGQ, and Cypher in Section~2.3 to offer practical context for graph query languages.
  \item Reworked Section~3 by introducing a consistent running example, streamlining the presentation of the core algebra operators, and adding concrete examples for selection, union, and concatenation.
  \item Improved Section~4 by adding a step-by-step worked example of recursive path construction via repeated concatenation, illustrating how short path fragments generate longer walks, and clarified the distinction between walks, trails, and acyclic paths using the example sequence.
  \item Extended Section~5 by adding a concrete translation of a GQL path pattern into algebraic operators and included an illustrative example showing additional expressiveness beyond GQL/SQL/PGQ through ordering within recursive path construction.
  \item Expanded Section~6 with a more critical discussion of the advantages and limitations of the proposed path algebra relative to relational algebra, automata-theoretic approaches, and semiring frameworks.
  \item Expanded Section~7 by discussing optimization opportunities, implementation challenges, recursive evaluation complexity, scalability, cost-based optimization, and practical adoption in graph database systems.
  \item Added additional literature references and citations throughout to better relate the presented concepts to the broader graph database and query language literature.
\end{itemize}